\section{Introduction}\label{sec:introduction}

\lipsum[1]~\cite{Scribe25}.\footnote{This is a footnote.}
\begin{equation*}
	fact(n) = \begin{cases}
		1                 & \text{if } n = 0 \\
		n \cdot fact(n-1) & \text{if } n > 0
	\end{cases}
\end{equation*}
The same recursive factorial can be written in C as shown in Listing~\ref{lst:fact}.
\begin{lstlisting}[language=C, caption={A recursive factorial in C, mirroring the definition above.}, label={lst:fact}]
int fact(int n) {
    if (n == 0) return 1;
    return n * fact(n - 1);
}
\end{lstlisting}

\begin{infobox}{Definition}
	The factorial of $n$ is the product of all positive integers up to $n$, with the base case $fact(0) = 1$.
\end{infobox}

\noindent Callout boxes come in a few colors out of the box: \infoboxinline{info}, \warnboxinline{warn}, \tipboxinline{tip}, and \noteboxinline{note}. Define your own with \texttt{\textbackslash scribedefinebox\{name\}\{color\}}.

\fbcomment{This is a comment in orange from Federico Bruzzone.}
\lipsum[1]
\jdcomment{This is another comment in blue from Jhon Doe.}
\lipsum[2]

